\chapter{Optimal Sampling-Based Planning: \rrtstar and Informed \rrtstar}
\label{ch:ch17}
% Function names for pseudocode are prefixed with "RrtStar" to stay unique
% across the book.
\SetKwFunction{RrtStarPlan}{RRTstar}
\SetKwFunction{RrtStarChooseParent}{ChooseParent}
\SetKwFunction{RrtStarRewire}{Rewire}
\SetKwFunction{RrtStarUpdate}{UpdateSubtree}
\SetKwFunction{RrtStarSample}{SampleInformed}
\SetKwFunction{RrtStarInformed}{InformedRRTstar}
\SetKwFunction{RrtStarRotation}{RotationToWorldFrame}

\begin{objectives}
\begin{itemize}
  \item Explain what \rrtstar adds to \rrt (\textsc{Near}, \textsc{ChooseParent}, \textsc{Rewire}), trace one iteration by hand, and say why each addition is needed for optimality.
  \item State the connection radius $r_n=\min\{\gamma_{\rrtstar}(\log n/n)^{1/d},\eta\}$ exactly, explain every symbol, compute an admissible $\gamma_{\rrtstar}$ for a given map, and explain why the radius must shrink slowly.
  \item State the asymptotic-optimality theorem with its assumptions, sketch the ball-covering proof, and say what the theorem does \emph{not} promise.
  \item Derive the informed set of \irrtstar (an ellipse in 2D, an ellipsoid in 3D), sample it uniformly with the scale--rotate--translate transformation, and predict when it helps.
  \item Implement \rrtstar and \irrtstar in \python in 2D and 3D, compare them with \rrt and with grid \astar, and use them as anytime planners while a drone is flying.
\end{itemize}
\end{objectives}

% ---------------------------------------------------------------------
\section{Why this matters for a drone swarm}
\label{sec:ch17-motivation}

The rapidly-exploring random tree of \cref{ch:ch16} solves a problem that
grid search cannot touch: it finds a collision-free route through a
continuous 3D volume with a few thousand samples, whatever the resolution a
grid would have needed. But look at the route it returns. It is a chain of
short, randomly oriented segments, it wanders, and on the map of
\cref{ch:ch16} it comes out somewhere between $20$ and $27$ units long when
the shortest route has length $16.9$. Smoothing helps, but it can only pull
the path taut inside the corridor the tree happened to find. Worse, more
samples do not help at all: the parent of every vertex is fixed at the moment
the vertex is created, so the first path that reaches the goal is, up to the
random order of insertion, the path you keep. A drone that flies a route
thirty percent longer than necessary spends thirty percent more battery on
every mission, and a swarm of them spends it thirty times over.

\rrtstar, introduced by Karaman and Frazzoli in 2011
\cite{karaman2011sampling}, repairs this with two extra steps per
iteration. Instead of attaching the new vertex to the nearest vertex, it
attaches it to the neighbour that gives the cheapest route from the start
(\textsc{ChooseParent}), and it then lets the new vertex become the parent
of any neighbour whose route it shortens (\textsc{Rewire}). The tree keeps
re-optimising itself, and its best path converges to the optimum as the
number of samples grows: \rrtstar is \emph{asymptotically optimal}. The price
is a logarithmic number of extra collision checks per iteration. \irrtstar
\cite{gammell2014informed} then adds the observation that, once a path of
cost $c_{\mathrm{best}}$ is known, every point that could improve it lies
inside an ellipsoid with the start and the goal as foci, and that sampling
this ellipsoid directly makes the convergence much faster in large free
spaces.

In the hybrid architecture of \cref{ch:ch24} these planners live in two
places. The global layer plans the swarm on a grid with \cbs or \ecbs
(\cref{ch:ch09,ch:ch10}); when a drone must leave that grid, because it
flies through a 3D volume too large or too finely structured to discretise,
\rrtstar plans the continuous segment, and it does so in \emph{anytime}
fashion: the first path is available after a few hundred samples, and the
tree keeps improving it while the drone is already flying. The replanning
layer can seed \irrtstar with the grid path it already has, so that the
sampler concentrates on the ellipsoid around that path from the first
iteration. This is the Week~8 topic of the study plan (\cref{ch:appA}):
its milestone is that you can choose between graph search and
sampling-based planning, and its coding exercise is \cref{exr:ch17-coding}.

The chapter proceeds as follows. \Cref{sec:ch17-intuition} gives the idea
and \cref{sec:ch17-problem} the definitions. \Cref{sec:ch17-algorithm}
presents the three new primitives, the pseudocode, the connection radius and
its $k$-nearest alternative, and traces one iteration by hand.
\Cref{sec:ch17-example} runs \rrtstar on the map of \cref{ch:ch16} and
compares it with \rrt and with grid \astar. \Cref{sec:ch17-properties} states
and sketches the asymptotic-optimality theorem and analyses the cost per
iteration. \Cref{sec:ch17-informed} derives \irrtstar,
\cref{sec:ch17-variants} surveys the family up to BIT\textsuperscript{*},
\cref{sec:ch17-implementation} covers the implementation and its pitfalls,
and \cref{sec:ch17-drone} places the planners in the drone system.

% ---------------------------------------------------------------------
\section{The idea in plain words}
\label{sec:ch17-intuition}

Think of the \rrt tree as a road network that is built one road at a time.
\rrt builds each new road from the nearest existing junction, whichever
junction that is. If the nearest junction happens to lie at the end of a
long detour, the new road inherits the detour, and so does every road built
from it later. Nothing in \rrt ever revisits this decision.

\rrtstar makes two changes, both local to the new vertex $x_{\mathrm{new}}$
and both restricted to a ball of radius $r_n$ around it
(\cref{fig:ch17-primitives}).
\begin{itemize}
  \item \textbf{Choose the parent.}\index{RRT*!ChooseParent} Among all
  vertices within $r_n$ of $x_{\mathrm{new}}$, connect to the one that gives
  the shortest route from the start, provided the connecting segment is
  collision-free. The nearest vertex is only the fallback.
  \item \textbf{Rewire the neighbours.}\index{RRT*!Rewire} For every vertex
  within $r_n$ of $x_{\mathrm{new}}$, check whether the route through
  $x_{\mathrm{new}}$ is shorter than its current route. If it is, and the
  segment is collision-free, make $x_{\mathrm{new}}$ its parent. The saving
  is passed on to all its descendants.
\end{itemize}
Because a vertex's parent can change long after the vertex was created,
the tree is no longer a record of the order in which space was explored; it
is an ever-improving \emph{shortest-path tree} over the sampled points, in
the same sense that Dijkstra's algorithm (\cref{ch:ch03}) maintains a
shortest-path tree over a graph. As samples accumulate, the sampled points
become dense, the shortest paths through them approach the shortest paths
through free space, and the cost of the best tree path to the goal
converges to the optimum.

The radius $r_n$ is the one delicate quantity. If it stays constant, the
number of neighbours grows linearly with the number of vertices $n$ and
every iteration does $\bigO{n}$ collision checks. If it shrinks too fast,
the balls become empty, neighbours are no longer found, and the tree stops
improving. The right schedule, $r_n\propto(\log n/n)^{1/d}$, shrinks as
slowly as it can while keeping the expected number of neighbours
logarithmic; it is the same scale at which a random geometric graph stays
connected \cite{penrose2003random}.

\begin{keyidea}
\rrtstar is \rrt with a shortest-path-tree discipline: every new vertex takes
the cheapest parent within radius $r_n$, and it becomes the parent of every
neighbour it can serve more cheaply. With $r_n\propto(\log n/n)^{1/d}$ and a
large enough constant, the best path converges to the optimum almost surely.
\irrtstar keeps the same tree but, once a path of cost $c_{\mathrm{best}}$ exists,
samples only the ellipsoid $\norm{x-x_{\mathrm{init}}}+\norm{x-x_{\mathrm{goal}}}\le c_{\mathrm{best}}$, the only region in which a better path can pass.
\end{keyidea}

\begin{figure}[tb]
  \centering
  \begin{subfigure}[b]{0.32\textwidth}\centering
    \inputfigure{ch17/near}
    \caption{\textsc{Near}}
  \end{subfigure}\hfill
  \begin{subfigure}[b]{0.32\textwidth}\centering
    \inputfigure{ch17/choose-parent}
    \caption{\textsc{ChooseParent}}
  \end{subfigure}\hfill
  \begin{subfigure}[b]{0.32\textwidth}\centering
    \inputfigure{ch17/rewire}
    \caption{\textsc{Rewire}}
  \end{subfigure}
  \caption{One \rrtstar iteration on the ten-vertex tree of
  \cref{ex:ch17-hand}. (a)~The ball of radius $r_n=2$ around
  $x_{\mathrm{new}}$ contains $v_3,v_4,v_6,v_7,v_8$; $v_4$ is the nearest
  vertex. (b)~The cost-to-come through each candidate: $v_6$ would be the
  cheapest parent ($4.58$) but the segment crosses the obstacle, so $v_7$
  wins with $5.26$, not the nearest $v_4$ ($5.69$). (c)~$v_8$ is cheaper
  through $x_{\mathrm{new}}$ ($6.25$ instead of $6.45$), so its edge from
  $v_4$ is cut and the saving reaches its child $v_9$.}
  \label{fig:ch17-primitives}
\end{figure}

% ---------------------------------------------------------------------
\section{Problem statement and notation}
\label{sec:ch17-problem}

We keep the setting of \cref{ch:ch16}. The planning space is a box
$\mathcal{X}\subset\R^d$ with $d\in\{2,3\}$ (a robot position, or a
configuration after the obstacles have been inflated by the drone radius as
in \cref{ch:ch02}), $\Xobs$ is its obstacle region and
$\Xfree=\mathcal{X}\setminus\Xobs$ the free space. The start is
$x_{\mathrm{init}}\in\Xfree$ and the goal $x_{\mathrm{goal}}\in\Xfree$.
A path is a continuous map $\sigma:[0,1]\to\mathcal{X}$; it is
\textbf{collision-free}\index{path!collision-free} if
$\sigma(s)\in\Xfree$ for all $s$, and its cost $c(\sigma)$ is its Euclidean
length. Throughout, $\mu(\cdot)$ denotes volume (Lebesgue measure),
$\disc{x}{r}$ the closed ball of radius $r$ around $x$, and
$\zeta_d=\mu(\disc{0}{1})$ the volume of the unit ball: $\zeta_2=\pi$,
$\zeta_3=4\pi/3$.

\begin{definition}[Optimal planning problem]
\label{def:ch17-problem}
Given $(\Xfree,x_{\mathrm{init}},x_{\mathrm{goal}})$, find a collision-free path
$\sigma^*$ from $x_{\mathrm{init}}$ to $x_{\mathrm{goal}}$ with
$c(\sigma^*)=c^*:=\inf\set{c(\sigma):\sigma\text{ collision-free from }x_{\mathrm{init}}\text{ to }x_{\mathrm{goal}}}$.
We write $c_{\min}=\norm{x_{\mathrm{goal}}-x_{\mathrm{init}}}$ for the
straight-line distance; every path has $c(\sigma)\ge c_{\min}$.
\end{definition}

Sampling-based planners can only find paths that have some room around
them, so the theory needs a notion of clearance.

\begin{definition}[Clearance and robust optimum]
\label{def:ch17-clearance}
A path $\sigma$ has \textbf{strong $\delta$-clearance}\index{clearance!strong}
if $\disc{\sigma(s)}{\delta}\subseteq\Xfree$ for all $s\in[0,1]$. An optimal
path $\sigma^*$ is a \textbf{robustly optimal solution}\index{robust optimality}
if there are collision-free paths $\sigma_n$ with strong
$\delta_n$-clearance, $\delta_n>0$, such that $c(\sigma_n)\to c^*$.
\end{definition}

The optimum of the map in \cref{ch:ch16} hugs the corners of the two walls,
so it has no clearance itself; but pushing it away from the corners by any
$\delta>0$ costs only $\bigO{\delta}$, so it is robustly optimal. A route
through a gap of width exactly zero is the kind of optimum that is
\emph{not} robust: no path with positive clearance approximates it, and no
sampling-based planner will find it.

\begin{definition}[Tree, cost-to-come, best cost]
\label{def:ch17-tree}
The planner maintains a tree $T=(V,E)$ rooted at $x_{\mathrm{init}}$ with
$n=\abs{V}$ vertices; $\mathrm{parent}(v)$ is the parent of $v\ne
x_{\mathrm{init}}$. The \textbf{cost-to-come}\index{cost-to-come}
$\mathrm{Cost}(v)$ is the length of the tree path from $x_{\mathrm{init}}$ to
$v$, so $\mathrm{Cost}(v)=\mathrm{Cost}(\mathrm{parent}(v))+\norm{v-\mathrm{parent}(v)}$
and $\mathrm{Cost}(x_{\mathrm{init}})=0$. The \textbf{best cost} after $n$
iterations is $c_{\mathrm{best}}=\mathrm{Cost}(x_{\mathrm{goal}})$ if
$x_{\mathrm{goal}}\in V$ and $c_{\mathrm{best}}=\infty$ otherwise.
\end{definition}

The primitives $\Nearest(V,x)$, $\Steer(x,y)$ (move from $x$ towards $y$ by
at most the step size $\eta$) and $\CollisionFree(x,y)$ (the straight
segment from $x$ to $y$ lies in $\Xfree$) are those of \cref{ch:ch16}.
\rrtstar adds one more.

\begin{definition}[Near and the connection radius]
\label{def:ch17-near}
$\Near(V,x,r)=\set{v\in V:\norm{v-x}\le r}$\index{RRT*!Near} is the set of
vertices within distance $r$ of $x$. \rrtstar calls it with the
\textbf{connection radius}\index{connection radius}
\begin{equation}
  r_n=\min\Bigl\{\gamma_{\rrtstar}\Bigl(\frac{\log n}{n}\Bigr)^{1/d},\ \eta\Bigr\},
  \qquad
  \gamma_{\rrtstar}>\gamma^*_{\rrtstar}:=2\Bigl(1+\frac1d\Bigr)^{1/d}\Bigl(\frac{\mu(\Xfree)}{\zeta_d}\Bigr)^{1/d},
  \label{eq:ch17-radius}
\end{equation}
where $n=\abs{V}$ is the current number of vertices, $d$ the dimension,
$\eta$ the step size of $\Steer$ and $\log$ the natural logarithm.
\end{definition}

\begin{definition}[Asymptotic optimality]
\label{def:ch17-ao}
Let $Y_n$ be the cost of the best solution in the tree after $n$
iterations ($Y_n=\infty$ if there is none). A planner is
\textbf{asymptotically optimal}\index{asymptotic optimality} if, for every
problem with a robustly optimal solution of cost $c^*$,
$\Prob\bigl(\lim_{n\to\infty}Y_n=c^*\bigr)=1$.
\end{definition}

Finally, the set that \irrtstar samples.

\begin{definition}[Informed set]
\label{def:ch17-informed}
For a current best cost $c_{\mathrm{best}}\ge c_{\min}$, the
\textbf{informed set}\index{informed set} is
\begin{equation}
  \mathcal{X}_{\hat f}(c_{\mathrm{best}})=\set{x\in\mathcal{X}:\ \norm{x-x_{\mathrm{init}}}+\norm{x-x_{\mathrm{goal}}}\le c_{\mathrm{best}}}.
  \label{eq:ch17-informed-set}
\end{equation}
\end{definition}

% ---------------------------------------------------------------------
\section{The algorithm}
\label{sec:ch17-algorithm}

\subsection{The three new primitives}
\label{sec:ch17-primitives}

An \rrtstar iteration begins exactly like an \rrt iteration: draw
$x_{\mathrm{rand}}$, find $x_{\mathrm{nearest}}=\Nearest(V,x_{\mathrm{rand}})$,
steer to $x_{\mathrm{new}}$, and give up on the sample if the segment from
$x_{\mathrm{nearest}}$ to $x_{\mathrm{new}}$ is blocked. If it is free,
\rrt would insert the edge $(x_{\mathrm{nearest}},x_{\mathrm{new}})$ and
stop. \rrtstar instead runs three steps, illustrated in
\cref{fig:ch17-primitives} on the small tree that \cref{ex:ch17-hand} works
through in numbers.

\paragraph{Near.} $X_{\mathrm{near}}=\Near(V,x_{\mathrm{new}},r_n)$ collects
every vertex within the connection radius of the new point. It is the
candidate set for both of the following steps; $x_{\mathrm{nearest}}$ is
added to it if the radius happens to exclude it, so that a parent always
exists.

\paragraph{ChooseParent.} For every candidate $v\in X_{\mathrm{near}}$
compute the cost-to-come that $x_{\mathrm{new}}$ would have through $v$,
namely $\mathrm{Cost}(v)+\norm{x_{\mathrm{new}}-v}$, and connect
$x_{\mathrm{new}}$ to the candidate with the smallest value whose segment
$\CollisionFree(v,x_{\mathrm{new}})$ holds. The implementation sorts the
candidates by that value and tests the segments in increasing order,
stopping at the first free one; most of the time the first test succeeds
and the step costs a single collision check.

\paragraph{Rewire.} For every candidate $v\in X_{\mathrm{near}}$ other than
the chosen parent, compute the cost $v$ would have through the new vertex,
$\mathrm{Cost}(x_{\mathrm{new}})+\norm{v-x_{\mathrm{new}}}$. If it is smaller
than $\mathrm{Cost}(v)$ and $\CollisionFree(x_{\mathrm{new}},v)$ holds,
replace the edge from $\mathrm{parent}(v)$ to $v$ by the edge from
$x_{\mathrm{new}}$ to $v$. The cost of $v$ drops, and so does the cost of
every descendant of $v$ by the same amount; \textsc{UpdateSubtree}
propagates the change through the subtree so that the invariant
$\mathrm{Cost}(u)=\mathrm{Cost}(\mathrm{parent}(u))+\norm{u-\mathrm{parent}(u)}$
holds again for every vertex.

\subsection{Pseudocode}
\label{sec:ch17-pseudocode}

\Cref{alg:ch17-rrtstar} is the main loop and \cref{alg:ch17-helpers} its
three helpers. The goal is handled as in \cref{ch:ch16} by goal
bias:\index{goal bias} with probability $p_{\mathrm{goal}}$ the sample is
$x_{\mathrm{goal}}$ itself, and since $\Steer$ returns the target when it is
within $\eta$, the goal point eventually becomes a vertex. From then on the
tree path to that vertex \emph{is} the current best path, its cost is
$c_{\mathrm{best}}$, and rewiring keeps shortening it; the bias is switched
off because the goal needs no second copy.

\begin{algorithm}[htb]
\caption{\rrtstar (Karaman and Frazzoli)}
\label{alg:ch17-rrtstar}
\KwIn{free space $\Xfree$ (via $\CollisionFree$), start $x_{\mathrm{init}}$, goal $x_{\mathrm{goal}}$, step $\eta$, constant $\gamma_{\rrtstar}$, goal bias $p_{\mathrm{goal}}$, budget $N$}
\KwOut{a tree $T$; the tree path to $x_{\mathrm{goal}}$ of cost $c_{\mathrm{best}}$, or \emph{failure}}
\Fn{\RrtStarPlan{$x_{\mathrm{init}}$, $x_{\mathrm{goal}}$, $N$}}{
  $V\gets\set{x_{\mathrm{init}}}$; $E\gets\emptyset$; $\mathrm{Cost}(x_{\mathrm{init}})\gets0$; $c_{\mathrm{best}}\gets\infty$\label{alg:ch17-rrtstar:init}\;
  \For{$i\gets1$ \KwTo $N$}{
    $x_{\mathrm{rand}}\gets x_{\mathrm{goal}}$ with probability $p_{\mathrm{goal}}$ if $x_{\mathrm{goal}}\notin V$, else a uniform sample of $\mathcal{X}$\label{alg:ch17-rrtstar:sample}\;
    $x_{\mathrm{nearest}}\gets\Nearest(V,x_{\mathrm{rand}})$; $x_{\mathrm{new}}\gets\Steer(x_{\mathrm{nearest}},x_{\mathrm{rand}})$\label{alg:ch17-rrtstar:steer}\;
    \If{$\CollisionFree(x_{\mathrm{nearest}},x_{\mathrm{new}})$\label{alg:ch17-rrtstar:free}}{
      $n\gets\abs{V}$; $r_n\gets\min\{\gamma_{\rrtstar}(\log n/n)^{1/d},\eta\}$\label{alg:ch17-rrtstar:radius}\;
      $X_{\mathrm{near}}\gets\Near(V,x_{\mathrm{new}},r_n)\cup\set{x_{\mathrm{nearest}}}$\label{alg:ch17-rrtstar:near}\;
      $x_{\mathrm{parent}}\gets$ \RrtStarChooseParent{$X_{\mathrm{near}}$, $x_{\mathrm{nearest}}$, $x_{\mathrm{new}}$}\label{alg:ch17-rrtstar:parent}\;
      $V\gets V\cup\set{x_{\mathrm{new}}}$; $E\gets E\cup\set{(x_{\mathrm{parent}},x_{\mathrm{new}})}$; $\mathrm{Cost}(x_{\mathrm{new}})\gets\mathrm{Cost}(x_{\mathrm{parent}})+\norm{x_{\mathrm{new}}-x_{\mathrm{parent}}}$\label{alg:ch17-rrtstar:add}\;
      \RrtStarRewire{$X_{\mathrm{near}}$, $x_{\mathrm{new}}$}\label{alg:ch17-rrtstar:rewire}\;
      \lIf{$x_{\mathrm{goal}}\in V$}{$c_{\mathrm{best}}\gets\mathrm{Cost}(x_{\mathrm{goal}})$\label{alg:ch17-rrtstar:best}}
    }
  }
  \lIf{$c_{\mathrm{best}}<\infty$}{\KwRet{the tree path from $x_{\mathrm{init}}$ to $x_{\mathrm{goal}}$}}
  \KwRet{failure}\;
}
\end{algorithm}

\begin{algorithm}[htb]
\caption{The helpers of \rrtstar}
\label{alg:ch17-helpers}
\Fn{\RrtStarChooseParent{$X_{\mathrm{near}}$, $x_{\mathrm{nearest}}$, $x_{\mathrm{new}}$}}{
  $x_{\min}\gets x_{\mathrm{nearest}}$; $c_{\min}\gets\mathrm{Cost}(x_{\mathrm{nearest}})+\norm{x_{\mathrm{new}}-x_{\mathrm{nearest}}}$\label{alg:ch17-helpers:default}\;
  \ForEach{$v\in X_{\mathrm{near}}$ in increasing order of $\mathrm{Cost}(v)+\norm{x_{\mathrm{new}}-v}$\label{alg:ch17-helpers:sort}}{
    \lIf{$\mathrm{Cost}(v)+\norm{x_{\mathrm{new}}-v}\ge c_{\min}$}{\KwBreak}
    \If{$\CollisionFree(v,x_{\mathrm{new}})$\label{alg:ch17-helpers:cpfree}}{
      $x_{\min}\gets v$; $c_{\min}\gets\mathrm{Cost}(v)+\norm{x_{\mathrm{new}}-v}$; \KwBreak\;
    }
  }
  \KwRet{$x_{\min}$}\;
}
\Fn{\RrtStarRewire{$X_{\mathrm{near}}$, $x_{\mathrm{new}}$}}{
  \ForEach{$v\in X_{\mathrm{near}}\setminus\set{\mathrm{parent}(x_{\mathrm{new}})}$}{
    $c\gets\mathrm{Cost}(x_{\mathrm{new}})+\norm{v-x_{\mathrm{new}}}$\label{alg:ch17-helpers:via}\;
    \If{$c<\mathrm{Cost}(v)$ \KwAnd $\CollisionFree(x_{\mathrm{new}},v)$\label{alg:ch17-helpers:rwfree}}{
      $E\gets E\setminus\set{(\mathrm{parent}(v),v)}\cup\set{(x_{\mathrm{new}},v)}$; $\mathrm{parent}(v)\gets x_{\mathrm{new}}$\label{alg:ch17-helpers:cut}\;
      \RrtStarUpdate{$v$}\label{alg:ch17-helpers:update}\;
    }
  }
}
\Fn{\RrtStarUpdate{$v$}}{
  $\mathrm{Cost}(v)\gets\mathrm{Cost}(\mathrm{parent}(v))+\norm{v-\mathrm{parent}(v)}$\;
  \lForEach{child $u$ of $v$}{\RrtStarUpdate{$u$}\tcp*[f]{depth-first over the subtree}}
}
\end{algorithm}

\paragraph{Walkthrough.} Lines~\ref{alg:ch17-rrtstar:sample}--\ref{alg:ch17-rrtstar:free}
are \rrt. Line~\ref{alg:ch17-rrtstar:radius} computes the radius from the
\emph{current} number of vertices, which is why the radius is recomputed in
every iteration, and line~\ref{alg:ch17-rrtstar:near} collects the
candidates. \textsc{ChooseParent} starts from the nearest vertex as the
default (line~\ref{alg:ch17-helpers:default}), whose segment is already known
to be free, and improves on it only with a candidate whose segment passes
the check in line~\ref{alg:ch17-helpers:cpfree}; the early \KwBreak after
the sorted loop means that no candidate is tested once it cannot beat the
best value found. The new vertex is inserted with its cost
(line~\ref{alg:ch17-rrtstar:add}) \emph{before} rewiring, because
\textsc{Rewire} reads $\mathrm{Cost}(x_{\mathrm{new}})$ in
line~\ref{alg:ch17-helpers:via}. Line~\ref{alg:ch17-helpers:rwfree} tests
the improvement first and the collision second, so segments are checked only
when they would be used; line~\ref{alg:ch17-helpers:cut} swaps the parent
and line~\ref{alg:ch17-helpers:update} pushes the new cost down the subtree.
Line~\ref{alg:ch17-rrtstar:best} reads $c_{\mathrm{best}}$ from the tree; it
can only decrease, since \textsc{Rewire} never increases any cost
(\cref{prop:ch17-invariant}).

\subsection{The connection radius}
\label{sec:ch17-radius}

Every symbol in \cref{eq:ch17-radius} has a job.
\begin{itemize}
  \item $n=\abs{V}$ is the number of vertices \emph{in the tree}, not the
  number of iterations; blocked samples add no vertex and do not shrink the
  radius.
  \item $d$ is the dimension. Volume scales like $r^d$, so a ball must have
  radius $\propto(\cdot)^{1/d}$ to contain a prescribed fraction of the space.
  \item $(\log n/n)^{1/d}$ is the schedule. A ball of radius $r_n$ has volume
  $\zeta_d r_n^d=\zeta_d\gamma^d\log n/n$, so with $n$ vertices spread over
  $\mu(\Xfree)$ the expected number of vertices in it is
  $n\cdot\zeta_d\gamma^d\log n/(n\,\mu(\Xfree))=\zeta_d\gamma^d\log n/\mu(\Xfree)$,
  which grows only logarithmically.
  \item $\gamma_{\rrtstar}$ is the constant that decides \emph{how many}
  neighbours that logarithm holds. The threshold $\gamma^*_{\rrtstar}$ makes
  the expected number of neighbours at least $2^d(1+1/d)\log n$; with this
  many, the ball-covering argument of \cref{sec:ch17-properties} goes through.
  \item $\mu(\Xfree)$ is the volume of the free space; the larger it is, the
  sparser the vertices, and the larger the radius must be to find neighbours.
  Any upper bound, such as $\mu(\mathcal{X})$, is a safe substitute because a
  larger $\gamma$ still satisfies the condition.
  \item $\zeta_d$ converts the radius into a volume.
  \item $\eta$ caps the radius at the step size, so that connections are
  never longer than the segments $\Steer$ produces. For small $n$ the
  logarithmic term is far larger than $\eta$ and the cap is active; the
  shrinking phase begins only when $\gamma(\log n/n)^{1/d}$ falls below
  $\eta$.
\end{itemize}

\begin{example}[The constants of the two maps]
\label{ex:ch17-gamma}
The map of \cref{ch:ch16} is the square $[0,10]^2$ with the walls
$[3,5]\times[0,6]$ and $[6,8]\times[4,10]$, so $\mu(\Xfree)=100-12-12=76$
and, with $d=2$ and $\zeta_2=\pi$,
$\gamma^*_{\rrtstar}=2\sqrt{1.5}\,\sqrt{76/\pi}=12.05$. We use
$\gamma_{\rrtstar}=13.3$, ten percent above the threshold. With $\eta=1$ the
cap is active until $13.3\sqrt{\log n/n}<1$, that is until about
$n\approx1\,650$ vertices; at $n=2\,193$ the radius is $0.79$. For the 3D
box field of \cref{ch:ch16}, $\mu(\mathcal{X})=1\,000$ is a safe upper bound
for $\mu(\Xfree)$ and gives $\gamma^*_{\rrtstar}=2(4/3)^{1/3}(1000/\zeta_3)^{1/3}=13.65$;
the code uses $15.0$.
\end{example}

\begin{figure}[tb]
  \centering
  \inputfigure{ch17/radius}
  \caption{Left: the connection radius of \cref{eq:ch17-radius} for the map
  of \cref{ex:ch17-gamma} ($\mu=76$, $d=2$, $\eta=1$). The cap $\eta$ is
  active while the tree is small; the dashed curve is the uncapped
  schedule. Right: the expected size of the \textsc{Near} set,
  $n\zeta_2r_n^2/\mu$. It grows linearly while the cap is active and then
  follows $\zeta_2\gamma^2\log n/\mu$ (dashed): a larger $\gamma$ buys more
  rewiring per iteration at a higher price per iteration.}
  \label{fig:ch17-radius}
\end{figure}

\paragraph{Why the radius must shrink, and shrink slowly.}\index{connection radius!shrinking}
The optimality proof needs two things at once. The tree must keep finding
neighbours everywhere along a near-optimal path, which pushes the radius up;
and the cost per iteration must stay bounded by a slowly growing function of
$n$, which pushes it down. Consider a ball of radius $r_n$ somewhere in the
free space. If $r_n$ shrank like $n^{-1/d}$ or faster, the expected number of
vertices in the ball, $\propto nr_n^d$, would stay bounded, and by the
Poisson approximation the ball would be empty with a probability bounded
away from zero, iteration after iteration. Somewhere along the optimal path
a ball would then stay empty for long stretches and the tree path would keep
its detour: this is the sense in which \emph{a radius that is too small
kills optimality} (\cref{sec:ch17-implementation}). The factor $\log n$ is
exactly what makes the expected count grow without bound, and the constant
in front decides how fast: with $\gamma$ above the threshold, the
probability that a given ball is empty decays like $n^{-c}$ with $c$ large
enough that, summed over all $n$ and over the $\bigO{n^{1/d}}$ balls needed
to cover the path, the sum is finite. On the other side, a radius that did
not shrink, such as $r_n=\eta$, would make the \textsc{Near} set grow like
$n$ and the total work like $n^2$. The schedule $(\log n/n)^{1/d}$ is the
slowest shrinking radius that keeps the work per iteration at
$\bigO{\log n}$, and it is not a coincidence that the same $\log n/n$ is the
threshold at which a random geometric graph on $n$ uniform points becomes
connected \cite{penrose2003random}: \rrtstar succeeds because the graph
whose edges are all pairs within $r_n$ is, with high probability, connected
along every path with clearance.

\subsection{The $k$-nearest variant}
\label{sec:ch17-knn}

Instead of a ball, \textsc{Near} may return the $k_n$ nearest
vertices,\index{RRT*!k-nearest@$k$-nearest} with
\begin{equation}
  k_n=\bigl\lceil k_{\mathrm{RRG}}\log n\bigr\rceil,\qquad k_{\mathrm{RRG}}>e\Bigl(1+\frac1d\Bigr),
  \label{eq:ch17-knn}
\end{equation}
where $e$ is Euler's number \cite{karaman2011sampling}. The threshold has
the practical merit of not depending on $\mu(\Xfree)$, which is often
unknown; $k_{\mathrm{RRG}}=2e\approx5.44$ is valid in every dimension, and
the code uses $1.1\,e(1+1/d)$. The trade is a different behaviour in
non-uniform regions: in a dense cluster the $k$ nearest vertices lie inside
a tiny ball and rewiring becomes very local, while in a sparse region the
$k$-th neighbour may be far away and its segment is likely to be blocked.
On the map of \cref{ex:ch17-gamma}, $k_n=33$ at $n=1\,438$ vertices, and
after $2\,000$ iterations the $k$-nearest version reaches a best cost of
$17.44$ against $17.45$ for the ball version with the same seed; on this map
the two are interchangeable.

\subsection{One iteration by hand}
\label{sec:ch17-hand}

\begin{example}[One \rrtstar iteration on a ten-vertex tree]
\label{ex:ch17-hand}
\Cref{fig:ch17-primitives} shows a tree with root $v_0=(0.5,0.5)$ in the
box $[0,6]\times[0,4.5]$ with one obstacle $[2.6,3.1]\times[1.9,2.45]$; the
first block of \cref{tab:ch17-hand} lists the vertices, their parents and
their costs. The sample is $x_{\mathrm{rand}}=(4.0,2.1)$ and $\eta=1$. The
nearest vertex is $v_4=(4.5,1.3)$ at distance $0.943<\eta$, so
$x_{\mathrm{new}}=x_{\mathrm{rand}}$, and the segment $v_4x_{\mathrm{new}}$ is
free. With $r_n=2$, $\Near$ returns $\set{v_3,v_4,v_6,v_7,v_8}$; the
distances to $x_{\mathrm{new}}$ are $1.900$, $0.943$, $1.924$, $1.221$ and
$0.985$. \textsc{ChooseParent} sorts the candidates by
$\mathrm{Cost}(v)+\norm{x_{\mathrm{new}}-v}$: $v_6$ ($2.654+1.924=4.578$),
$v_7$ ($5.264$), $v_3$ ($5.440$), $v_4$ ($5.692$), $v_8$ ($7.436$). The
segment from $v_6$ crosses the obstacle and is rejected; the segment from
$v_7$ is free, so $x_{\mathrm{new}}$ becomes $v_{10}$ with parent $v_7$ and
cost $5.264$, $0.43$ less than through the nearest vertex. \textsc{Rewire}
then tests the remaining candidates: $v_3$, $v_4$ and $v_6$ are cheaper as
they are, but $v_8$ would cost $5.264+0.985=6.249<6.451$, and its segment is
free. Its parent changes from $v_4$ to $v_{10}$, and \textsc{UpdateSubtree}
lowers its child $v_9$ from $7.583$ to $7.380$ (second block of
\cref{tab:ch17-hand}). Every number is produced by \code{tiny\_example()} in
\code{code/ch17\_rrt\_star.py}.
\end{example}

\begin{table}[tb]
  \centering\small
  \caption{Trace of \cref{ex:ch17-hand}: the tree before and after one
  iteration of \cref{alg:ch17-rrtstar}. Costs are rounded to three
  decimals; changed entries are in bold.}
  \label{tab:ch17-hand}
  \begin{tabular}{@{}lccrccr@{}}
  \toprule
  & \multicolumn{3}{c}{before} & \multicolumn{3}{c}{after}\\
  \cmidrule(lr){2-4}\cmidrule(lr){5-7}
  Vertex & position & parent & cost & parent & cost & comment\\
  \midrule
  $v_0$ & $(0.5,0.5)$ & -- & $0$ & -- & $0$ & root\\
  $v_1$ & $(1.6,0.2)$ & $v_0$ & $1.140$ & $v_0$ & $1.140$ & \\
  $v_2$ & $(2.8,0.2)$ & $v_1$ & $2.340$ & $v_1$ & $2.340$ & \\
  $v_3$ & $(4.0,0.2)$ & $v_2$ & $3.540$ & $v_2$ & $3.540$ & near; $7.164$ via $v_{10}$, no\\
  $v_4$ & $(4.5,1.3)$ & $v_3$ & $4.748$ & $v_3$ & $4.748$ & nearest; $6.207$ via $v_{10}$, no\\
  $v_5$ & $(0.9,1.7)$ & $v_0$ & $1.265$ & $v_0$ & $1.265$ & \\
  $v_6$ & $(2.1,2.4)$ & $v_5$ & $2.654$ & $v_5$ & $2.654$ & near; cheapest parent but blocked\\
  $v_7$ & $(3.3,3.1)$ & $v_6$ & $4.043$ & $v_6$ & $4.043$ & near; chosen parent\\
  $v_8$ & $(4.4,3.0)$ & $v_4$ & $6.451$ & $\mathbf{v_{10}}$ & $\mathbf{6.249}$ & near; rewired\\
  $v_9$ & $(5.2,3.8)$ & $v_8$ & $7.583$ & $v_8$ & $\mathbf{7.380}$ & updated with its parent\\
  $v_{10}$ & $(4.0,2.1)$ & -- & -- & $\mathbf{v_7}$ & $\mathbf{5.264}$ & $x_{\mathrm{new}}$\\
  \bottomrule
  \end{tabular}
\end{table}

% ---------------------------------------------------------------------
\section{A worked example}
\label{sec:ch17-example}

\begin{example}[\rrtstar on the map of \cref{ch:ch16}]
\label{ex:ch17-map}
The map is the square $[0,10]^2$ with the walls $[3,5]\times[0,6]$ and
$[6,8]\times[4,10]$, the start is $(1,1)$ and the goal $(9,9)$; the robot is
a point and segments are checked at spacing $\delta=0.05$ with the
conservative inflation of \cref{ch:ch16}. The only route passes over the
first wall and under the second, and the optimum, found by Dijkstra on the
visibility graph of the wall corners (pushed out by the checker's
inflation), has length $c^*=16.909$. We run \cref{alg:ch17-rrtstar} with
$\eta=1$, $\gamma_{\rrtstar}=13.3$, $p_{\mathrm{goal}}=0.05$ and seed~$1$ for
$3\,000$ iterations, recording the tree and the best path at iterations
$200$, $1\,000$ and $3\,000$ (\cref{fig:ch17-snapshots}) and
$c_{\mathrm{best}}$ at every iteration (\cref{tab:ch17-convergence}).
\end{example}

\begin{figure}[tb]
  \centering
  \inputfigure{ch17/snapshots}
  \caption{\rrtstar on \cref{ex:ch17-map} after (a)~$200$, (b)~$1\,000$ and
  (c)~$3\,000$ iterations, with the current best path in blue. Compare the
  tree with the \rrt tree of \cref{ch:ch16} on the same map: the edges point
  radially away from the start, because every vertex hangs from the parent
  that gives it the shortest route, and the best path straightens as
  rewiring continues, from $18.21$ at $1\,000$ iterations to $17.32$ at
  $3\,000$ against the optimum $16.91$.}
  \label{fig:ch17-snapshots}
\end{figure}

\begin{table}[tb]
  \centering\small
  \caption{Best cost against iterations for \cref{ex:ch17-map} (seed~$1$),
  with the number of vertices and the connection radius at that moment. The
  radius stays at the cap $\eta=1$ until about $1\,650$ vertices and then
  follows $13.3\sqrt{\log n/n}$. The last column is plain \rrt with the same
  seed and sampling: its first path is its only path.}
  \label{tab:ch17-convergence}
  \begin{tabular}{@{}rrrrr@{}}
  \toprule
  Iteration & $\abs{V}$ & $r_n$ & $c_{\mathrm{best}}$ (\rrtstar) & $c_{\mathrm{best}}$ (\rrt)\\
  \midrule
  $200$ & $103$ & $1.000$ & -- & --\\
  $351$ & $204$ & $1.000$ & $21.003$ & $25.721$\\
  $500$ & $311$ & $1.000$ & $19.726$ & $25.721$\\
  $1\,000$ & $680$ & $1.000$ & $18.212$ & $25.721$\\
  $1\,500$ & $1\,058$ & $1.000$ & $17.711$ & $25.721$\\
  $2\,000$ & $1\,438$ & $0.946$ & $17.452$ & $25.721$\\
  $3\,000$ & $2\,193$ & $0.788$ & $17.324$ & $25.721$\\
  \midrule
  optimum $c^*$ & & & $16.909$ & $16.909$\\
  \bottomrule
  \end{tabular}
\end{table}

The first path appears at iteration $351$, when the goal point becomes the
$204$th vertex, and costs $21.003$; \rrt with the same seed reaches the goal
at the same iteration, because both planners insert the same vertices, but
its path costs $25.721$ and never changes. From then on \rrtstar improves at
every rewiring that touches the goal's branch: $19.73$ at $500$ iterations,
$18.21$ at $1\,000$, $17.32$ at $3\,000$, that is $2.5\%$ above the optimum.
The improvement slows down as the tree densifies, as the theory predicts:
the remaining gap is made of many small corner-cutting steps, each of which
needs a sample in a small region. Over twenty seeds the final cost of
\rrtstar is $17.33\pm0.08$ (range $17.18$ to $17.49$) while that of \rrt is
$23.51\pm1.67$ (range $20.75$ to $26.74$): the random tree of \rrt makes the
path quality a lottery, \rrtstar makes it a matter of budget.

The Week~8 plan also asks for a comparison with grid \astar. On the same map
discretised at $h=0.1$ with 8-connected moves and no corner cutting, \astar
(\cref{ch:ch04}) returns a path of length $18.29$: it is optimal on the grid,
but the octile detours of the grid make it $8\%$ longer than the continuous
optimum. String pulling (the greedy pruning of \cref{ch:ch16}) shortens it to
$17.23$, and no further, because the grid keeps the path one cell away from
every corner. \rrtstar reaches the quality of the pulled grid path at about
$2\,500$ iterations and keeps going; in the same experiment it needed
$7\,997$ segment checks for $3\,000$ iterations, $2.7$ times the $3\,000$ of
\rrt, with a \textsc{Near} set of $23.6$ vertices on average.
